\documentclass[11pt,a4paper]{article}

% Packages for styling and formatting
\usepackage[utf8]{inputenc}
\usepackage[T1]{fontenc}
\usepackage{lmodern}
\usepackage[margin=1in]{geometry}
\usepackage{xcolor}
\usepackage{hyperref}
\usepackage{graphicx}
\usepackage{booktabs}
\usepackage{tabularx}
\usepackage{titlesec}
\usepackage{parskip}
\usepackage{enumitem}
\usepackage{fancyhdr}
\usepackage{tcolorbox}
\usepackage{microtype}
\usepackage{listings}

% Color definitions
\definecolor{primary}{RGB}{3, 105, 161} % Dark blue
\definecolor{secondary}{RGB}{217, 119, 6} % Orange/Amber
\definecolor{darkgray}{RGB}{24, 24, 27}
\definecolor{lightgray}{RGB}{244, 244, 245}
\definecolor{getcolor}{RGB}{16, 185, 129} % Green
\definecolor{postcolor}{RGB}{59, 130, 246} % Blue
\definecolor{putcolor}{RGB}{245, 158, 11} % Orange
\definecolor{deletecolor}{RGB}{239, 68, 68} % Red
\definecolor{codebg}{RGB}{248, 250, 252}

% Hyperref setup
\hypersetup{
    colorlinks=true,
    linkcolor=primary,
    filecolor=primary,      
    urlcolor=primary,
    pdftitle={API Documentation: Music Appreciation API},
    pdfpagemode=FullScreen,
}

% Section styling
\titleformat{\section}
  {\normalfont\Large\bfseries\color{primary}}
  {\thesection}{1em}{}
  [\vspace{-0.5ex}\color{secondary}\titlerule]

\titleformat{\subsection}
  {\normalfont\large\bfseries\color{darkgray}}
  {\thesubsection}{1em}{}

% Header and Footer
\pagestyle{fancy}
\fancyhf{}
\fancyhead[L]{\textcolor{darkgray}{\small XJCO3011 Coursework 1}}
\fancyhead[R]{\textcolor{darkgray}{\small API Documentation}}
\fancyfoot[C]{\thepage}
\renewcommand{\headrulewidth}{0.4pt}
\renewcommand{\footrulewidth}{0.4pt}

% Listings setup for JSON
\lstdefinelanguage{json}{
    basicstyle=\ttfamily\small,
    numbers=none,
    stepnumber=1,
    numbersep=8pt,
    showstringspaces=false,
    breaklines=true,
    frame=single,
    backgroundcolor=\color{codebg},
    rulecolor=\color{lightgray},
    stringstyle=\color{primary},
    keywordstyle=\color{secondary}\bfseries,
}

% Custom commands for HTTP methods
\newcommand{\methodGET}{\colorbox{getcolor}{\textcolor{white}{\textbf{GET}}}}
\newcommand{\methodPOST}{\colorbox{postcolor}{\textcolor{white}{\textbf{POST}}}}
\newcommand{\methodPUT}{\colorbox{putcolor}{\textcolor{white}{\textbf{PUT}}}}
\newcommand{\methodDELETE}{\colorbox{deletecolor}{\textcolor{white}{\textbf{DELETE}}}}

\begin{document}

% Title Page
\begin{titlepage}
    \centering
    \vspace*{2cm}
    
    {\Huge\bfseries\color{primary} Music Appreciation and Discovery API\par}
    \vspace{1cm}
    {\huge\bfseries API Documentation\par}
    \vspace{1.5cm}
    
    {\Large\bfseries XJCO3011 - Web Services and Web Data\par}
    \vspace{0.5cm}
    {\large Coursework 1 - Individual Web Services API Development Project\par}
    
    \vspace{2cm}
    
    {\Large\bfseries Zeyu Shen\par}
    \vspace{0.5cm}
    {\large \today\par}
    
    \vfill
    
    \begin{tcolorbox}[colback=white,colframe=secondary,boxrule=1pt,arc=4pt,width=0.8\textwidth,center]
    \centering
    \textbf{Base URL:} \texttt{https://weidademiaoxiao.pythonanywhere.com} \\
    \textbf{Authentication:} API Key required for write operations (\texttt{X-API-Key}).
    \end{tcolorbox}
    
    \vspace{1cm}
\end{titlepage}

\newpage
\tableofcontents
\newpage

\section{Overview}

The Music Appreciation and Discovery API is a RESTful web service designed to facilitate the exploration of music tracks, genres, and user-generated content such as reviews, tags, and collections. It provides a robust set of endpoints for data retrieval, CRUD operations, and analytical insights.

This document outlines the available endpoints, their parameters, and expected responses.

\section{Endpoints}

\subsection{1. General}

\subsubsection*{\methodGET \texttt{ /}}
Returns a welcome message and links to the interactive documentation.
\begin{itemize}[label=\textcolor{primary}{\textbullet}, noitemsep]
    \item \textbf{Response (200 OK):}
\begin{lstlisting}[language=json]
{
  "message": "Music Appreciation and Discovery API is running.",
  "docs": "/docs",
  "redoc": "/redoc",
  "version": "0.3.0"
}
\end{lstlisting}
\end{itemize}

\subsubsection*{\methodGET \texttt{ /health}}
Returns a simple health-check response to verify the API is operational.
\begin{itemize}[label=\textcolor{primary}{\textbullet}, noitemsep]
    \item \textbf{Response (200 OK):}
\begin{lstlisting}[language=json]
{
  "status": "ok"
}
\end{lstlisting}
\end{itemize}

\subsection{2. Genres}

\subsubsection*{\methodGET \texttt{ /genres}}
Retrieves a list of all available music genres, ordered alphabetically.
\begin{itemize}[label=\textcolor{primary}{\textbullet}, noitemsep]
    \item \textbf{Response (200 OK):} Array of Genre objects.
\begin{lstlisting}[language=json]
[
  {
    "id": 1,
    "name": "Ambient",
    "description": "Atmospheric and meditative soundscapes..."
  }
]
\end{lstlisting}
\end{itemize}

\subsubsection*{\methodGET \texttt{ /genres/\{genre\_id\}}}
Retrieves a specific genre by its ID.
\begin{itemize}[label=\textcolor{primary}{\textbullet}, noitemsep]
    \item \textbf{Parameters:} \texttt{genre\_id} (path, integer)
    \item \textbf{Response (200 OK):} Genre object.
    \item \textbf{Response (404 Not Found):} If the genre does not exist.
\end{itemize}

\subsection{3. Tracks}

\subsubsection*{\methodGET \texttt{ /tracks}}
Retrieves a list of tracks with optional filtering, sorting, and pagination.
\begin{itemize}[label=\textcolor{primary}{\textbullet}, noitemsep]
    \item \textbf{Query Parameters:}
    \begin{itemize}[label=--, noitemsep]
        \item \texttt{title} (string, optional): Filter by track title (partial match).
        \item \texttt{artist} (string, optional): Filter by artist name (partial match).
        \item \texttt{genre} (string, optional): Filter by genre name (partial match).
        \item \texttt{mood} (string, optional): Filter by mood keyword (partial match).
        \item \texttt{sort\_by} (string, default="title"): Sort field (\texttt{title}, \texttt{artist}, \texttt{rating}, \texttt{play\_count}).
        \item \texttt{order} (string, default="asc"): Sort order (\texttt{asc} or \texttt{desc}).
        \item \texttt{limit} (integer, default=20, max=100): Maximum number of results.
        \item \texttt{offset} (integer, default=0): Number of results to skip.
    \end{itemize}
    \item \textbf{Response (200 OK):} Array of Track objects.
\end{itemize}

\subsubsection*{\methodGET \texttt{ /tracks/\{track\_id\}}}
Retrieves a specific track by its ID.
\begin{itemize}[label=\textcolor{primary}{\textbullet}, noitemsep]
    \item \textbf{Parameters:} \texttt{track\_id} (path, integer)
    \item \textbf{Response (200 OK):} Track object including nested genre details.
    \item \textbf{Response (404 Not Found):} If the track does not exist.
\end{itemize}

\subsection{4. Reviews}

\subsubsection*{\methodPOST \texttt{ /reviews}}
Creates a new review for a track. Requires API Key.
\begin{itemize}[label=\textcolor{primary}{\textbullet}, noitemsep]
    \item \textbf{Headers:} \texttt{X-API-Key: <your-api-key>}
    \item \textbf{Request Body:}
\begin{lstlisting}[language=json]
{
  "track_id": 1,
  "reviewer_name": "John Doe",
  "rating": 5,
  "comment": "An absolute masterpiece."
}
\end{lstlisting}
    \item \textbf{Response (201 Created):} The created Review object.
    \item \textbf{Response (404 Not Found):} If the specified track does not exist.
    \item \textbf{Response (422 Unprocessable Entity):} If validation fails (e.g., rating not between 1 and 5).
\end{itemize}

\subsubsection*{\methodGET \texttt{ /reviews}}
Retrieves a list of reviews with optional filtering and pagination.
\begin{itemize}[label=\textcolor{primary}{\textbullet}, noitemsep]
    \item \textbf{Query Parameters:}
    \begin{itemize}[label=--, noitemsep]
        \item \texttt{track\_id} (integer, optional): Filter reviews by track ID.
        \item \texttt{reviewer\_name} (string, optional): Filter by reviewer name (partial match).
        \item \texttt{min\_rating} (integer, optional): Minimum rating filter (1-5).
        \item \texttt{limit} (integer, default=20, max=100): Maximum number of results.
        \item \texttt{offset} (integer, default=0): Number of results to skip.
    \end{itemize}
    \item \textbf{Response (200 OK):} Array of Review objects.
\end{itemize}

\subsubsection*{\methodGET \texttt{ /reviews/\{review\_id\}}}
Retrieves a specific review by its ID.
\begin{itemize}[label=\textcolor{primary}{\textbullet}, noitemsep]
    \item \textbf{Parameters:} \texttt{review\_id} (path, integer)
    \item \textbf{Response (200 OK):} Review object.
    \item \textbf{Response (404 Not Found):} If the review does not exist.
\end{itemize}

\subsubsection*{\methodPUT \texttt{ /reviews/\{review\_id\}}}
Updates an existing review. Supports partial updates. Requires API Key.
\begin{itemize}[label=\textcolor{primary}{\textbullet}, noitemsep]
    \item \textbf{Headers:} \texttt{X-API-Key: <your-api-key>}
    \item \textbf{Parameters:} \texttt{review\_id} (path, integer)
    \item \textbf{Request Body:} (All fields optional)
\begin{lstlisting}[language=json]
{
  "rating": 4,
  "comment": "Updated comment text."
}
\end{lstlisting}
    \item \textbf{Response (200 OK):} The updated Review object.
    \item \textbf{Response (404 Not Found):} If the review does not exist.
\end{itemize}

\subsubsection*{\methodDELETE \texttt{ /reviews/\{review\_id\}}}
Deletes a review by its ID. Requires API Key.
\begin{itemize}[label=\textcolor{primary}{\textbullet}, noitemsep]
    \item \textbf{Headers:} \texttt{X-API-Key: <your-api-key>}
    \item \textbf{Parameters:} \texttt{review\_id} (path, integer)
    \item \textbf{Response (204 No Content):} Successful deletion.
    \item \textbf{Response (404 Not Found):} If the review does not exist.
\end{itemize}

\subsection{5. Tags}

\subsubsection*{\methodPOST \texttt{ /tags}}
Creates a new user tag for a track. Requires API Key.
\begin{itemize}[label=\textcolor{primary}{\textbullet}, noitemsep]
    \item \textbf{Headers:} \texttt{X-API-Key: <your-api-key>}
    \item \textbf{Request Body:}
\begin{lstlisting}[language=json]
{
  "track_id": 1,
  "tag_name": "masterpiece",
  "created_by": "Jane Doe"
}
\end{lstlisting}
    \item \textbf{Response (201 Created):} The created Tag object.
    \item \textbf{Response (404 Not Found):} If the track does not exist.
    \item \textbf{Response (409 Conflict):} If the tag already exists for this track.
\end{itemize}

\subsubsection*{\methodGET \texttt{ /tags}}
Retrieves a list of user tags, optionally filtered by track.
\begin{itemize}[label=\textcolor{primary}{\textbullet}, noitemsep]
    \item \textbf{Query Parameters:} \texttt{track\_id} (integer, optional)
    \item \textbf{Response (200 OK):} Array of Tag objects.
\end{itemize}

\subsubsection*{\methodDELETE \texttt{ /tags/\{tag\_id\}}}
Deletes a user tag by its ID. Requires API Key.
\begin{itemize}[label=\textcolor{primary}{\textbullet}, noitemsep]
    \item \textbf{Headers:} \texttt{X-API-Key: <your-api-key>}
    \item \textbf{Parameters:} \texttt{tag\_id} (path, integer)
    \item \textbf{Response (204 No Content):} Successful deletion.
    \item \textbf{Response (404 Not Found):} If the tag does not exist.
\end{itemize}

\subsection{6. Collections}

\subsubsection*{\methodPOST \texttt{ /collections}}
Creates a new track collection. Requires API Key.
\begin{itemize}[label=\textcolor{primary}{\textbullet}, noitemsep]
    \item \textbf{Headers:} \texttt{X-API-Key: <your-api-key>}
    \item \textbf{Request Body:}
\begin{lstlisting}[language=json]
{
  "name": "Late Night Study",
  "description": "Calm tracks for focused work.",
  "created_by": "Alice"
}
\end{lstlisting}
    \item \textbf{Response (201 Created):} The created Collection object.
    \item \textbf{Response (409 Conflict):} If a collection with the same name already exists.
\end{itemize}

\subsubsection*{\methodGET \texttt{ /collections}}
Retrieves a list of all collections, ordered by creation date.
\begin{itemize}[label=\textcolor{primary}{\textbullet}, noitemsep]
    \item \textbf{Response (200 OK):} Array of Collection objects.
\end{itemize}

\subsubsection*{\methodGET \texttt{ /collections/\{collection\_id\}}}
Retrieves a specific collection along with its track items.
\begin{itemize}[label=\textcolor{primary}{\textbullet}, noitemsep]
    \item \textbf{Parameters:} \texttt{collection\_id} (path, integer)
    \item \textbf{Response (200 OK):} Collection object including an array of \texttt{items}.
    \item \textbf{Response (404 Not Found):} If the collection does not exist.
\end{itemize}

\subsubsection*{\methodPOST \texttt{ /collections/\{collection\_id\}/items}}
Adds a track to an existing collection. Requires API Key.
\begin{itemize}[label=\textcolor{primary}{\textbullet}, noitemsep]
    \item \textbf{Headers:} \texttt{X-API-Key: <your-api-key>}
    \item \textbf{Parameters:} \texttt{collection\_id} (path, integer)
    \item \textbf{Request Body:}
\begin{lstlisting}[language=json]
{
  "track_id": 1,
  "note": "Great opening track."
}
\end{lstlisting}
    \item \textbf{Response (201 Created):} Confirmation message and item details.
    \item \textbf{Response (404 Not Found):} If the collection or track does not exist.
    \item \textbf{Response (409 Conflict):} If the track is already in the collection.
\end{itemize}

\subsubsection*{\methodDELETE \texttt{ /collections/\{collection\_id\}}}
Deletes a collection and all its items. Requires API Key.
\begin{itemize}[label=\textcolor{primary}{\textbullet}, noitemsep]
    \item \textbf{Headers:} \texttt{X-API-Key: <your-api-key>}
    \item \textbf{Parameters:} \texttt{collection\_id} (path, integer)
    \item \textbf{Response (204 No Content):} Successful deletion.
    \item \textbf{Response (404 Not Found):} If the collection does not exist.
\end{itemize}

\subsubsection*{\methodDELETE \texttt{ /collections/\{collection\_id\}/items/\{item\_id\}}}
Removes a specific track item from a collection. Requires API Key.
\begin{itemize}[label=\textcolor{primary}{\textbullet}, noitemsep]
    \item \textbf{Headers:} \texttt{X-API-Key: <your-api-key>}
    \item \textbf{Parameters:} 
    \begin{itemize}[label=--, noitemsep]
        \item \texttt{collection\_id} (path, integer)
        \item \texttt{item\_id} (path, integer)
    \end{itemize}
    \item \textbf{Response (204 No Content):} Successful deletion.
    \item \textbf{Response (404 Not Found):} If the collection item does not exist.
\end{itemize}

\subsection{7. Analytics}

\subsubsection*{\methodGET \texttt{ /analytics/top-rated-tracks}}
Returns the highest-rated tracks that have at least one review.
\begin{itemize}[label=\textcolor{primary}{\textbullet}, noitemsep]
    \item \textbf{Query Parameters:} \texttt{limit} (integer, default=5, max=20)
    \item \textbf{Response (200 OK):} Array of track summary objects.
\end{itemize}

\subsubsection*{\methodGET \texttt{ /analytics/genre-summary}}
Returns a summary of each genre, including track count and average rating.
\begin{itemize}[label=\textcolor{primary}{\textbullet}, noitemsep]
    \item \textbf{Response (200 OK):} Array of genre summary objects.
\end{itemize}

\subsubsection*{\methodGET \texttt{ /analytics/top-tags}}
Returns the most frequently used tags across all tracks.
\begin{itemize}[label=\textcolor{primary}{\textbullet}, noitemsep]
    \item \textbf{Query Parameters:} \texttt{limit} (integer, default=10, max=50)
    \item \textbf{Response (200 OK):} Array of tag summary objects.
\end{itemize}

\subsubsection*{\methodGET \texttt{ /analytics/mood-distribution}}
Returns the distribution of moods across all tracks.
\begin{itemize}[label=\textcolor{primary}{\textbullet}, noitemsep]
    \item \textbf{Response (200 OK):} Array of mood summary objects.
\end{itemize}

\subsubsection*{\methodGET \texttt{ /analytics/review-activity}}
Returns overall review activity summary and per-reviewer breakdown.
\begin{itemize}[label=\textcolor{primary}{\textbullet}, noitemsep]
    \item \textbf{Response (200 OK):} Object containing \texttt{total\_reviews}, \texttt{average\_rating}, and an array of \texttt{reviewers}.
\end{itemize}

\section{Data Models}

\subsection{Track}
\begin{lstlisting}[language=json]
{
  "id": 1,
  "title": "Clair de Lune",
  "artist_name": "Claude Debussy",
  "album_title": "Suite bergamasque",
  "duration_seconds": 300,
  "mood": "dreamy",
  "play_count": 1500,
  "average_rating": 4.5,
  "genre": {
    "id": 1,
    "name": "Classical",
    "description": "..."
  }
}
\end{lstlisting}

\subsection{Review}
\begin{lstlisting}[language=json]
{
  "id": 1,
  "track_id": 1,
  "reviewer_name": "John Doe",
  "rating": 5,
  "comment": "An absolute masterpiece.",
  "created_at": "2023-10-27T10:00:00Z",
  "updated_at": "2023-10-27T10:00:00Z"
}
\end{lstlisting}

\end{document}
